% NS-53. Exact identities, proved implications, and an explicitly open input.
% No finite numerical approximation is offered as convergence evidence.
\subsection{NS-53: residualized integer-dilation blocks}
\label{sec:ns53-nb-blocks}

Work in the real Hilbert space $H=L^2((0,\infty),dx)$, with
$\chi=\mathbf 1_{(0,1)}$ and $\rho_n(x)=\{1/(nx)\}$ for integers $n\geq1$.
Real coefficients lose nothing for this real target; the complex version
uses conjugate transposes. Put
\[
 V_N=\operatorname{span}\{\rho_1,\ldots,\rho_N\},\qquad
 r_N=\chi-P_N\chi,\qquad E_N=\|r_N\|^2=d_N^2,
\]
where $P_N$ is the orthogonal projection onto $V_N$.
The strong Nyman--Beurling criterion is the established equivalence
$E_N\to0\Longleftrightarrow\mathrm{RH}$; see B\'aez-Duarte,
\href{https://arxiv.org/abs/math/0202141}{Theorem~1.1}.
No particular upper rate is required. Burnol's
\href{https://doi.org/10.1006/aima.2001.2066}{lower-bound theorem}
does not require a matching upper bound. This checkpoint proves no
convergence statement and makes no G2 or RH claim.

\begin{proposition}[Exact arithmetic entries and finite independence]
\label{prop:ns53-nb-entries}
Let $G_{mn}=\langle\rho_m,\rho_n\rangle$, $b_n=\langle\chi,\rho_n\rangle$,
and $\kappa=\log(2\pi)-\gamma$, where $\gamma$ is Euler's constant.
Then
\begin{align}
 G_{mn}&=\int_0^\infty\{y/m\}\{y/n\}\,\frac{dy}{y^2},
 & b_n&=\frac{\log n+1-\gamma}{n},
 &G_{nn}&=\frac{\kappa}{n}.\label{eq:ns53-nb-entries}
\end{align}
In particular, writing $a_j=\lfloor j/m\rfloor$ and
$d_j=\lfloor j/n\rfloor$, the following is an exact convergent series:
\begin{equation}
 G_{mn}=\frac1{mn}+\sum_{j=1}^{\infty}
 \left[\frac1{mn}-\left(\frac{d_j}{m}+\frac{a_j}{n}\right)
 \log\left(1+\frac1j\right)+\frac{a_jd_j}{j(j+1)}\right].
 \label{eq:ns53-nb-gram-series}
\end{equation}
Each bracket is a nonnegative cell integral. Every Gram matrix belonging
to finitely many distinct positive integers is positive definite.
Moreover $E_N>0$ for every finite $N$.
\end{proposition}

\begin{proof}
Substitute $y=1/x$ and integrate on $[j,j+1)$; both floor functions
are constant on that interval. This gives the integral and series formulas.
The change $y=nt$ gives the diagonal scaling and
\[
 b_n=\frac1n\left(\log n+\int_1^\infty\frac{\{t\}}{t^2}\,dt\right).
\]
The integral truncated at an integer $K$ is $1+\log K-H_K$.
Similarly, integrating $\{t\}^2/t^2$ on the first $M+1$ unit cells gives
\[
 2M+2-H_{M+1}-2M\log(M+1)+2\log(M!).
\]
Stirling's formula gives its limit $\kappa$. For finite independence,
a vanishing linear combination of $\{y/n\}$ is piecewise affine.
At the smallest index with nonzero coefficient its jump cannot be
cancelled by a smaller nonzero coefficient, a contradiction. If
$\chi=\sum_{n\leq N}c_n\rho_n$, comparison of the jumps at every positive
integer in the $y$ variable gives
$\sum_{n\mid j}c_n=-\mathbf 1_{j=1}$, with $c_n=0$ for $n>N$.
M\"obius inversion would force $c_n=-\mu(n)$ for every $n$, contradicted
by a prime greater than $N$. A finite-dimensional span is closed, so
$E_N>0$.
\end{proof}

\begin{proposition}[Exact block improvement, including all cross terms]
\label{prop:ns53-nb-block-gain}
Fix integers $1\leq N<M$. Let $G$ be the $N$ by $N$ old Gram matrix,
$C=(G_{mn})_{1\leq m\leq N,\,N<n\leq M}$, and
$D=(G_{mn})_{N<m,n\leq M}$.
Write $b=(b_m)_{m\leq N}$, $b_B=(b_n)_{N<n\leq M}$, and define
\begin{equation}
 c=G^{-1}b,\qquad h=b_B-C^{\mathsf T}c,\qquad
 S=D-C^{\mathsf T}G^{-1}C.
 \label{eq:ns53-nb-schur}
\end{equation}
Then $S$ is positive definite, and
\begin{equation}
 E_N-E_M=h^{\mathsf T}S^{-1}h
 =\sup_{w\ne0}\frac{|w^{\mathsf T}h|^2}{w^{\mathsf T}Sw}.
 \label{eq:ns53-nb-gain}
\end{equation}
The gain is zero if and only if $h=0$; strict finite improvement is not
asserted for every block. The coefficients for the enlarged optimum are
$(c-G^{-1}CS^{-1}h,S^{-1}h)$. In addition,
\begin{equation}
 E_N-E_M\ \geq\
 \frac{\sum_{n=N+1}^M|b_n-\sum_{m\leq N}c_mG_{mn}|^2}
 {\kappa\sum_{n=N+1}^M1/n}
 \ \geq\
 \frac{\|h\|_2^2}{\kappa\log(M/N)}.
 \label{eq:ns53-nb-trace-gain}
\end{equation}
These are finite, unconditional identities and inequalities, not estimates
of the numerator as $N$ grows.
\end{proposition}

\begin{proof}
Set $z_n=(I-P_N)\rho_n$ for $N<n\leq M$. Its Gram matrix is $S$,
and $\langle r_N,z_n\rangle=h_n$. Finite independence shows that these
$z_n$ are independent. The orthogonal decomposition
$V_M=V_N\mathbin{\oplus}\operatorname{span}\{z_n\}$ gives
\eqref{eq:ns53-nb-gain} and the coefficient update. Since
$0<S\leq D$, the largest eigenvalue of $S$ is at most
$\operatorname{tr}D=\kappa\sum_{n=N+1}^M1/n$. Therefore
$h^{\mathsf T}S^{-1}h\geq\|h\|_2^2/\operatorname{tr}D$.
The decreasing function $1/x$ gives
$\sum_{n=N+1}^M1/n\leq\log(M/N)$.
\end{proof}

\begin{remark}[Conditioning is retained]
Replacing $S$ by $D$ in the exact identity, or using $b_B$ instead of $h$,
would discard the old--new interaction. A small eigenvalue of $S$ can
make the exact gain much larger than the trace lower bound; it causes no
missing inverse error in the analytic lower bound. For a proposed real
block vector $w$, its exact gain after optimizing its scalar coefficient
is $|w^{\mathsf T}h|^2/(w^{\mathsf T}Sw)$. If
$\sum_{n=N+1}^M w_n/n=0$, the unprojected combination
$\sum w_n\rho_n$ vanishes on $x>1/(N+1)$. Its residualized combination
need not have this support, and its norm is still $w^{\mathsf T}Sw$.
\end{remark}

\begin{proposition}[Exact divisor cells and the raw M\"obius-prefix penalty]
\label{prop:ns53-nb-divisor-cells}
For arbitrary real coefficients $c_1,\ldots,c_N$, put
\[
 A=\sum_{n\leq N}\frac{c_n}{n},\qquad
 B_j=1+\sum_{n\leq N}c_n\left\lfloor\frac jn\right\rfloor
     =1+\sum_{k=1}^j\ \sum_{\substack{n\mid k\\n\leq N}}c_n.
\]
For $r=\chi-\sum c_n\rho_n$, one has $r(1/y)=-Ay$ on $(0,1)$
and $r(1/y)=B_j-Ay$ on $[j,j+1)$, $j\geq1$. Consequently
\begin{equation}
 \|r\|^2=A^2+\sum_{j=1}^{\infty}
 \left[A^2-2AB_j\log\left(1+\frac1j\right)
          +\frac{B_j^2}{j(j+1)}\right].
 \label{eq:ns53-nb-divisor-energy}
\end{equation}
The brackets are nonnegative cell integrals; they are not to be split
into three separately summed series. For the raw choice $c_n=-\mu(n)$,
\begin{equation}
 \left\|\chi+\sum_{n\leq N}\mu(n)\rho_n\right\|^2
 \geq (N+1)\left|\sum_{n\leq N}\frac{\mu(n)}n\right|^2.
 \label{eq:ns53-nb-mobius-penalty}
\end{equation}
\end{proposition}

\begin{proof}
Expand each fractional part. The energy identity is the integral of the
squared affine function on each unit cell. For the M\"obius choice,
$B_j=0$ for $1\leq j\leq N$, by
$\sum_{n\mid k}\mu(n)=\mathbf 1_{k=1}$. The interval $(0,N+1)$ alone
then contributes the displayed lower bound.
\end{proof}

\begin{remark}[What the arithmetic lemma does and does not supply]
Equation~\eqref{eq:ns53-nb-mobius-penalty} gives an explicit necessary
condition for this particular unsmoothed coefficient choice to converge:
$\sqrt N\sum_{n\leq N}\mu(n)/n\to0$. It is not a lower bound for the
optimal $E_N$. The divergence of the raw natural approximation is already
recorded in B\'aez-Duarte's Introduction; no novelty is claimed for its
failure. Smoothing or adaptive coefficients change the problem and must
be evaluated through the full residual and Gram matrix.
\end{remark}

\begin{proposition}[A sufficient block-correlation input at any rate]
\label{prop:ns53-nb-any-rate}
Fix an integer $q\geq2$, $N_j=q^j$, and real numbers
$0\leq\alpha_j<1$ for $j\geq j_0$, with $\sum_j\alpha_j=\infty$.
Use the actual optimal coefficients $c^{(N_j)}=G_{N_j}^{-1}b_{N_j}$.
Suppose the following \emph{residual block-correlation input} holds:
\begin{equation}
 \sum_{n=N_j+1}^{qN_j}
 \left|\frac{\log n+1-\gamma}{n}
       -\sum_{m\leq N_j}c_m^{(N_j)}G_{mn}\right|^2
 \geq\kappa\log q\,\alpha_j E_{N_j}.
 \tag{RBC}\label{eq:ns53-nb-rbc}
\end{equation}
Then $E_N\to0$. For example, the hypothetical bound (RBC) with
$\alpha_j=1/(2(j+1))$ would give
$E_{q^j}=O(j^{-1/2})$, hence $E_N=O((\log N)^{-1/2})$.
This slower-than-optimal rate would already meet the established criterion.
\end{proposition}

\begin{proof}
Equation~\eqref{eq:ns53-nb-trace-gain} gives
$E_{N_{j+1}}\leq(1-\alpha_j)E_{N_j}$. Iteration and
$1-t\leq e^{-t}$ prove convergence. Monotonicity supplies intermediate
$N$. For the example, summing $1/(2(j+1))$ yields the stated rate.
\end{proof}

\paragraph{Named open input and disposition.}
(RBC) is \emph{unproved}; it is a specific sufficient arithmetic estimate,
not asserted equivalent to RH or known to follow from RH. It may be
stronger than needed because the trace bound can lose the small
Gram eigenvalues that amplify the exact gain. The weaker exact target is
a nonsummable sequence of relative gains
$h_j^{\mathsf T}S_j^{-1}h_j/E_{N_j}$, which is equivalent to convergence
and is not itself a new theorem. The arithmetic work still missing is
a lower estimate for the corrected correlations, or for an explicitly
constructed block quotient, for the \emph{actual optimal residuals}.
Entry formulas, a finite collection of positive gains, and a finite list of distances
do not supply it. No such asymptotic gain estimate is proved here.
The missing step is the asymptotic gain estimate; this establishes
neither failure of (RBC) nor failure of the approximation object. No new Weil window, tail metric, G2 or RH
claim is introduced.
